\begin{Definition}
\label{def:essential-self-adjointness-and-core}
若对称算子 $(\mathcal{D}(A),A)$ 的闭包是自伴随的, 则称它是\textbf{本质自伴随的}. 若 $(\mathcal{D}(A),A)$ 是闭的, 则定义域的子集 $D\subset\mathcal{D}(A)$ 在满足 $\overline{(D,A)}=(\mathcal{D}(A),A)$ 时称为该算子的一个\textbf{核}; 即算子在 $D$ 上的限制的闭包为原算子.
\end{Definition}

本质自伴随性的最简单判别方法涉及称为亏指数的量.

\begin{Definition}
\label{def:operator-deficiency-indices}
设 $(\mathcal{D}(A),A)$ 为从 $H$ 到 $H$ 的稠密定义算子, 令 $\gamma^+$ 为 $\mathcal{R}(A-iI)^\perp$ 的维数, $\gamma^-$ 为 $\mathcal{R}(A+iI)^\perp$ 的维数. 数对 $(\gamma^+,\gamma^-)$ 称为该算子的\textbf{亏指数}.
\end{Definition}

\begin{Theorem}
\label{thm:essential-self-adjointness-zero-deficiency}
算子 $(\mathcal{D}(A),A)$ 是本质自伴随的, 当且仅当它的两个亏指数均为 $0$.
\end{Theorem}

证明留给读者 (Problem~\ref{ex:essential-self-adjointness-deficiency-proof}).

\subsubsection{伴随算子与谱理论}\label{subsec:adjoint-operators-spectral-theory}

通过研究算子的伴随, 可以对其谱有所了解. 第一个结果将算子的压缩谱与其 Hilbert 伴随的点谱联系起来.

\begin{Theorem}
\label{thm:adjoint-deficiency-conjugate-eigenvalue}
设 $(\mathcal{D}(A),A)$ 为从 $H$ 到 $H$ 的稠密定义算子. 复数 $\lambda$ 的亏数为 $m$, 当且仅当 $\overline\lambda$ 是 $A^*$ 的重数为 $m$ 的特征值.
\end{Theorem}

其证明是 Theorem~\ref{thm:range-complement-adjoint-nullspace} 的简单应用, 留给读者 (Problem~\ref{ex:adjoint-deficiency-eigenvalue-proof}).

关于对称算子的谱, 可以得到许多结论.

\begin{Theorem}
\label{thm:symmetric-operator-spectral-properties}
设 $(\mathcal{D}(A),A)$ 为从 $H$ 到 $H$ 的稠密定义算子. 若 $(\mathcal{D}(A),A)$ 是对称的, 则:\begin{enumerate}\item 对每个 $x\in\mathcal{D}(A)$, $(Ax,x)$ 都是实数.\item $A$ 的所有特征值都是实数.\item $A$ 的不同特征值所对应的特征向量彼此正交.\item $A$ 的连续谱是实的.\end{enumerate}
\end{Theorem}

\begingroup
\postdisplaypenalty=10000
\renewcommand{\endProof}{\par\nobreak\hfill$\blacksquare$\par\vspace{0.1cm}}
\begin{Proof}
为证明 1, 注意到对于每个 $x\in\mathcal{D}(A)$ 均有
\begin{equation}\label{eq:symmetric-operator-quadratic-form-real}
(Ax,x)=(x,Ax)=\overline{(Ax,x)}.
\end{equation}
为证明 2, 注意到若 $\lambda$ 是特征值且 $x\in\mathcal{N}(A_\lambda)$ 为对应的特征函数, 则利用它们是实数这一事实, 有
\begin{equation}\label{eq:symmetric-operator-eigenvalue-rayleigh-quotient}
\lambda=\dfrac{(x,Ax)}{\|x\|^2}.
\end{equation}
因而由第 1 部分, $\lambda$ 为实数.

为证明 3, 设 $\lambda_1$ 和 $\lambda_2$ 为 $A$ 的特征值, $x_1\in\mathcal{N}(A_{\lambda_1})$ 和 $x_2\in\mathcal{N}(A_{\lambda_2})$ 为对应的特征向量. 那么
\begin{equation}\label{eq:symmetric-operator-eigenvector-pairing}
\lambda_1(x_1,x_2)=(Ax_1,x_2)=(x_1,Ax_2)=\lambda_2(x_1,x_2).
\end{equation}
因而要么 $\lambda_1=\lambda_2$, 要么 $(x_1,x_2)=0$.

为证明 4, 令 $\lambda=\gamma+i\mu$, 其中 $\gamma$ 和 $\mu$ 为实数. 那么利用 $A$ 的对称性, 可以证明对于 $x\in\mathcal{D}(A)$ 有
\begin{equation}\label{eq:symmetric-operator-nonreal-parameter-lower-bound}
\|(A-\lambda I)x\|^2=\|Ax-\gamma x\|^2+\mu^2\|x\|^2\geq\mu^2\|x\|
\end{equation}
若 $|\mu|>0$, 则 $A-\lambda I$ 有下界. 因而由 Problem~\ref{ex:surjective-bounded-below-inverse}, $R_\lambda(A)$ 存在且有界. 证明完毕.
\end{Proof}
\endgroup

若算子是自伴随的, 还能得到更多结论.

\begin{Theorem}
\label{thm:self-adjoint-spectrum-real-no-compression}
设 $(\mathcal{D}(A),A)$ 为从 $H$ 到 $H$ 的稠密定义算子. 若 $(\mathcal{D}(A),A)$ 是自伴随的, 则每个虚部非零的 $\lambda\in\mathbb{C}$ 都属于 $A$ 的预解集. 此外, 压缩谱为空集.
\end{Theorem}

\begingroup
\postdisplaypenalty=10000
\renewcommand{\endProof}{\par\nobreak\hfill$\blacksquare$\par\vspace{0.1cm}}
\begin{Proof}
首先注意, Theorem~\ref{thm:symmetric-operator-spectral-properties} 表明 $A$ 的连续谱是实的, 且 $A$ 的所有特征值都是实数. 其次, Theorem~\ref{thm:adjoint-deficiency-conjugate-eigenvalue} 表明, 若 $\lambda$ 的亏数非零, 则 $\overline\lambda$ 是 $A(=A^*)$ 的特征值. 因而 $\lambda$ 必为实数, 且必属于点谱而非压缩谱.
\end{Proof}
\endgroup
